\chapter{Metabolism, Energy Gradients, and the Ecosystem}
\label{sec:metabolism}
% ===================================================================

\section{The External Gradient as Broken Symmetry}

The replicating fixed point of Section~\ref{sec:replication} is a
self-sustaining process, but it is not yet a \emph{living} one in
the full sense.
It copies itself, but copying costs account resources, and a closed
system with no external input will eventually exhaust its account and
fall silent.
Life requires not just replication but \emph{metabolism}: a
sustained coupling to an external energy source that replenishes
the account faster than computation depletes it.

The key physical observation, due to Eric Smith, is that the
Earth is not a closed system.
The Sun presents an enormous influx of energy --- a directed,
asymmetric gradient between the high-energy photon flux arriving
from the Sun and the low-energy thermal radiation leaving to space.
This gradient is not merely a source of replenishment; it is a
\emph{broken symmetry} in the weight map of the system.

Recall from Part~I that the reversibility construction
$\GSLT \to \GSLT^\dagger$ requires $\bwd_r(\phi) = \overline{\fwd_r(\phi)}$
for every rule $r$ --- the backward amplitude is the complex
conjugate of the forward amplitude.
This is the CPT condition: the system is time-symmetric.
The solar gradient violates this condition for a specific class
of rules.

\begin{definition}[Gradient-Coupled Rule]
A rewrite rule $r$ is \emph{gradient-coupled} if its forward
amplitude is boosted by an external energy flux $\Phi$:
\[
  w^+_r(\phi,\, \Phi) \;=\; w^{+(0)}_r(\phi) \cdot e^{\beta\Phi}
\]
where $w^{+(0)}_r$ is the amplitude in the absence of the
gradient and $\beta > 0$ is a coupling constant.
The corresponding backward amplitude remains
$w^-_r(\phi) = \overline{w^{+(0)}_r(\phi)}$, unaffected by the
gradient.
\end{definition}

\begin{remark}[CPT Violation as the Signature of Life]
A gradient-coupled rule has $|\fwd_r| > |\bwd_r|$: forward
transitions are more probable than backward ones.
The system is driven away from equilibrium.
This violation of the CPT condition of $\GSLT^\dagger$ is not a
defect of the framework but a \emph{signal}: it marks the
transitions through which the external gradient does work on the
system.
A living agent is precisely one that maintains a sustained
CPT-violating coupling to an external gradient.
An agent in thermal equilibrium with its environment --- one for
which $|\fwd_r| = |\bwd_r|$ for all rules --- is, in this
framework, dead.
\end{remark}

\section{Producers: Direct Gradient Coupling}

The simplest living agents are those that couple directly to the
external gradient.
In biological terms these are autotrophs --- plants and
photosynthetic bacteria.
In the GSLT framework they are agents whose metabolic rules have
the gradient $\Phi$ as part of their minimal context.

\begin{definition}[Producer]
A \emph{producer} is an agent $A$ that has at least one
gradient-coupled rewrite rule $r$ whose minimal context
$K_\Phi$ encodes the external gradient directly:
\[
  \langle A,\, \vect{A} \rangle
  \;\xrightarrow{K_\Phi}\;
  \langle A',\, \vect{A} + \delta\vect{A} \rangle
\]
where $\delta\vect{A} > \vect{0}$ componentwise.
The account grows: the producer converts gradient energy into
stored computational resource.
\end{definition}

\begin{remark}
The minimal context $K_\Phi$ for a producer rule is the encoding
of the gradient itself --- photons, in the biological case.
This is an instance of the open synchronization tree structure
from Part~I: the producer's metabolic transition is an edge in
$\mathcal{ST}_o(A)$ labeled by the gradient context.
The gradient is the environment; the producer is the program
that responds to it.
In the Rho calculus, $K_\Phi[-] = \Phi!(- ) \mid [-]$ where
$\Phi$ is the channel on which gradient quanta arrive.
A producer listens on $\Phi$ and credits its account on each
received quantum.
\end{remark}

\section{Consumers and the Food Chain as Account Flow}

Once producers have stored account resources in their processes,
other agents can couple to them rather than to the raw gradient.
These are consumers --- heterotrophs in biological terms.

\begin{definition}[Consumer]
A \emph{consumer} is an agent $B$ that has gradient-coupled
rules whose minimal context encodes not the external gradient
directly but the \emph{presence of a producer} (or of another
consumer at a lower trophic level):
\[
  \langle B,\, \vect{A}_B \rangle \;\inter\;
  \langle A,\, \vect{A}_A \rangle
  \;\rewrite\;
  \langle B',\, \vect{A}_B + \delta\vect{A} \rangle \;\inter\;
  \langle A',\, \vect{A}_A - \delta\vect{A} \rangle
\]
Account is transferred from $A$ to $B$ via interaction.
The total account is conserved across the interaction (no external
gradient is tapped); the distribution changes.
\end{definition}

\begin{remark}[Metabolism as Redistribution]
Smith's observation is precisely captured by this definition.
The Sun does not directly power animals: it powers plants, which
store account, which animals then redistribute.
Metabolism at the ecosystem level is not replenishment but
\emph{redistribution} of account through a population, with
producers as the sole point of contact with the external gradient.

The food chain is therefore a \emph{directed graph of account
transfers}: an edge from $A$ to $B$ means $B$ consumes $A$,
transferring stored account from $A$'s process to $B$'s.
The roots of this graph (the nodes with no incoming edges) are
the producers, coupled to the gradient.
The leaves are the apex consumers.
Decomposers close the cycle by returning stored account to
forms accessible to producers.
\end{remark}

\section{The Ecosystem as Account-Flow Network}

\begin{definition}[Ecosystem]
An \emph{ecosystem} is a compact coherence cluster
$\mathcal{C}$ (in the sense of Part~II) together with:
\begin{enumerate}[label=(\roman*)]
  \item a gradient $\Phi$ providing an external account source;
  \item a set of producers coupling $\mathcal{C}$ to $\Phi$;
  \item a directed account-flow graph $\mathcal{F}$ on
        $\mathcal{C}$ encoding who consumes whom.
\end{enumerate}
The ecosystem is \emph{viable} if the total account inflow from
$\Phi$ through the producers exceeds the total account dissipated
by the population's rewrites over any sufficiently long time
horizon.
\end{definition}

\begin{remark}[The Nerve as Metabolic Map]
The nerve $\mathcal{N}$ of Part~II --- the simplicial complex of
overlapping coherence clusters --- now has a metabolic
interpretation.
Each vertex of $\mathcal{N}$ is a coherence cluster; each edge
is a shared overlap.
When the clusters are populations of living agents, the edges of
$\mathcal{N}$ are also the channels of metabolic exchange: the
overlapping agents are those that can coherently communicate
\emph{and} transfer account resources.
The nerve is simultaneously the epistemic map (who can reach
consensus with whom) and the metabolic map (who can feed whom).
The two structures are not coincidentally aligned: coherent
communication is the precondition for reliable account transfer,
since a corrupted message about food availability is as
dangerous as no message at all.
\end{remark}

\section{Trophic Level as Causal Depth}

\begin{proposition}[Trophic Level as Open Synchronization Depth]
\label{prop:trophic}
The \emph{trophic level} of an agent $B$ in an ecosystem is
the minimum path length in the account-flow graph $\mathcal{F}$
from a producer to $B$.
This is precisely the minimum causal depth $d^o_{\min}$ in the
open synchronization tree of $B$ with respect to the gradient
context $K_\Phi$:
\[
  \mathrm{TL}(B) \;=\; d^o_{\min}(K_\Phi, B).
\]
\end{proposition}

\begin{remark}
Plants have trophic level 1: one context step separates them
from the gradient.
Herbivores have trophic level 2: their metabolic coupling to
the gradient passes through the plant context.
Carnivores that eat herbivores have trophic level 3; and so on.
Each trophic level is one additional layer of context nesting
in the open synchronization tree --- one more environmental
intermediary between the agent and the ultimate energy source.
The assembly index and the trophic level are thus both instances
of the same underlying quantity: causal depth in an open
synchronization tree, measured from different roots.
Assembly index is depth from the elementary terms.
Trophic level is depth from the gradient.
\end{remark}

\section{Viability of the Replicating Fixed Point}

The phase transition to life, as described in
Section~\ref{sec:phase}, produces a replicating fixed point
centered in a compact coherence cluster.
We can now state more precisely what it means for this fixed
point to be \emph{viable} in the long run.

\begin{conjecture}[Metabolic Viability of the Fixed Point]
\label{conj:viable}
A replicating fixed point $\mathbf{Y}(F)$ is \emph{metabolically
viable} if there exists a gradient $\Phi$ and a producer
sub-process $F_\Phi \subseteq F$ such that the account inflow
from $\Phi$ via $F_\Phi$ exceeds the account cost of replication
and maintenance:
\[
  \mathbb{E}\bigl[\delta\vect{A}_{\text{catabolic}}\bigr]
  \;>\;
  \mathbb{E}\bigl[\vect{c}_{\text{anabolic}}\bigr]
\]
where the expectations are over the path integral of the process.
Only metabolically viable fixed points persist after the phase
transition; the others exhaust their accounts and go silent.
\end{conjecture}

\begin{remark}[Selection for Metabolic Efficiency]
Conjecture~\ref{conj:viable} implies that the phase transition
does not merely select for replication but for
\emph{metabolically sustained replication}.
The Darwinian competition after the transition is not only
about who replicates fastest but about who maintains the best
ratio of catabolic yield to anabolic construction cost.
High assembly index processes win this competition when their
catalytic sophistication yields proportionally higher catabolic
returns --- which is precisely why complexity increases after
the phase transition.
The joint rise of $\mathrm{AI}$ and $\mathrm{CN}$ is driven not
just by selection for efficient replication but for efficient
\emph{metabolism}: the ability to extract more account from the
gradient per unit of construction cost.
\end{remark}

% ===================================================================
